% !TeX program = xelatex
% ============================================================================
%  第三部分  组合数学
%  第 1 章  计数原理与排列组合
%  第 2 章  二项式定理与组合恒等式
%  第 3 章  容斥原理
%  第 4 章  递推关系与经典计数数列
%  第 5 章  抽屉原理与极端原理
%  第 6 章  格路计数与反射原理
%  第 7 章  组合构造与综合例题
% ============================================================================

\chapter{计数原理与排列组合}
\begin{chapterintro}[章节导读]
组合数学的第一课不是公式，而是“把要数的东西说清楚”。同一道题，把对象理解成“有序”还是“无序”、把接收者理解成“有身份”还是“无身份”，答案会差一个阶乘。本章先把加法原理、乘法原理、排列、组合与插板法这条主线走通，并用“数的是什么”作为每一步的自查口号。
\end{chapterintro}

\begin{tip}[组合数学的总口诀]
\textbf{先说清数的是什么，再选工具。}具体做法：拿出一张纸，写下“一个被计数的对象由哪些独立的选择决定”“这些选择有没有顺序”“有没有重复计数”。这三问答完，方法基本就唯一了。本书第二部分数论里反复强调的“把整数看成指数向量”，本质也是同一件事。
\end{tip}

\section{加法原理与乘法原理}
\begin{theorem}[加法原理]
完成一件事有 $n$ 类互斥的办法，第 $i$ 类有 $a_i$ 种，则总方法数为
\[
S=a_1+a_2+\cdots+a_n.
\]
\end{theorem}
\begin{theorem}[乘法原理]
完成一件事分 $n$ 个步骤，第 $i$ 步有 $a_i$ 种方法，则总方法数为
\[
S=a_1a_2\cdots a_n.
\]
\end{theorem}
两条原理本身平淡无奇，难点全在\textbf{分类是否互斥、分步是否独立}。

\begin{pitfall}[分类不互斥]
把“$1$ 到 $20$ 中能被 $2$ 或被 $3$ 整除的数的个数”算成 $10+6=16$ 就错了：$6,12,18$ 被数了两次，正确答案是 $10+6-3=13$。这正是第三部分第 3 章容斥原理的出发点。\textbf{每次分类后问一句：有没有对象同时落进两类？}
\end{pitfall}

\section{排列与组合}
\begin{definition}[排列数]
从 $n$ 个不同元素中取出 $m$ 个（$m\le n$）排成一列，排列数为
\[
A_n^m=n(n-1)(n-2)\cdots(n-m+1)=\frac{n!}{(n-m)!}.
\]
特别地，$m=n$ 时叫全排列，$A_n^n=n!$。
\end{definition}
\begin{definition}[组合数]
从 $n$ 个不同元素中取出 $m$ 个组成一个集合（不计顺序），组合数为
\[
\binom nm=\frac{A_n^m}{m!}=\frac{n!}{m!(n-m)!}.
\]
读作“$n$ 选 $m$”。规定 $m>n$ 时 $\binom nm=0$。
\end{definition}

\begin{proof}[组合数公式的推导]
$n$ 个人选 $m$ 个排队有 $A_n^m$ 种；若不计顺序，则同一个人数集合被它的 $m!$ 种排列重复计算了 $m!$ 次，故 $\binom nm=A_n^m/m!$。
\end{proof}

\begin{tip}[排列与组合的唯一区别]
\textbf{顺序算不算不同。}排队、排名、放有编号的位置——顺序算不同，用排列；选人、选子集、分组——顺序不算，用组合。做题时先把这句话写在旁边，能避免一半的错误。
\end{tip}

\section{组合数的基本性质}
\begin{theorem}[组合数的五条基本性质]
\begin{align}
\binom nm&=\binom n{n-m}&&\text{（对称性）},\\
\binom nk&=\frac nk\binom{n-1}{k-1}&&\text{（递推，}k\ge1\text{）},\\
\binom nm&=\binom{n-1}{m}+\binom{n-1}{m-1}&&\text{（杨辉三角）},\\
\sum_{k=0}^{n}\binom nk&=2^n&&\text{（二项式定理，}a=b=1\text{）},\\
\sum_{k=0}^{n}(-1)^k\binom nk&=\iva{n=0}&&\text{（二项式定理，}a=1,b=-1\text{）}.
\end{align}
\end{theorem}

\begin{proof}[杨辉恒等式的组合证明]
要证 $\dbinom{n}{m}=\dbinom{n-1}{m}+\dbinom{n-1}{m-1}$，只需分别数“从 $n$ 个人里选 $m$ 个”的方法数。把某一个人（记为甲）单独拿出来：
\begin{itemize}
\item 不含甲的选法：从剩下 $n-1$ 个人里选 $m$ 个，共 $\dbinom{n-1}{m}$ 种；
\item 含甲的选法：甲已选定，再从剩下 $n-1$ 个人里选 $m-1$ 个，共 $\dbinom{n-1}{m-1}$ 种。
\end{itemize}
两类互斥且穷尽，相加即得。这个证明是\textbf{组合证明}的样板：不化简代数式，而是对同一个量给出两种不同的数法。
\end{proof}

\begin{teach}[代数证明与组合证明都要会]
杨辉恒等式也可以纯代数地证：
\[
\binom{n-1}{m}+\binom{n-1}{m-1}
=\frac{(n-1)!}{m!(n-1-m)!}+\frac{(n-1)!}{(m-1)!(n-m)!}
=\frac{n!}{m!(n-m)!}=\binom nm.
\]
两种证明都要求学生写一遍。代数证明训练恒等变形，组合证明训练“把同一个量数两遍”——后者是第三部分第 3 章容斥原理与第 5 章抽屉原理的思想源头。
\end{teach}

\section{插板法}
\begin{theorem}[插板法]
把 $n$ 个\textbf{相同}的元素分成 $k$ 组：
\begin{itemize}
\item 每组至少一个：$\dbinom{n-1}{k-1}$ 种；
\item 每组允许为空：$\dbinom{n+k-1}{k-1}=\dbinom{n+k-1}{n}$ 种；
\item 第 $i$ 组至少 $a_i$ 个（$\sum a_i\le n$）：$\dbinom{n-\sum a_i+k-1}{k-1}$ 种。
\end{itemize}
\end{theorem}
\begin{proof}
$n$ 个相同元素排成一列，中间有 $n-1$ 个空隙。每组至少一个，就是在这 $n-1$ 个空隙里插入 $k-1$ 块板子，每种插法对应一种分组，共 $\dbinom{n-1}{k-1}$ 种。

允许为空时不能直接插（多块板子可插入同一空隙），改用“借元素”：先借 $k$ 个元素，在 $n+k$ 个元素形成的 $n+k-1$ 个空隙里插 $k-1$ 块板子，保证每组至少一个，插完再把借来的 $k$ 个元素从各组中拿走。元素相同，所以这是一个双射，答案仍为 $\dbinom{n+k-1}{k-1}$。

不同下界时令 $x_i'=x_i-a_i$，化为非负整数解，套上一种情形。
\end{proof}

\begin{pitfall}[插板法的前提是“元素相同”]
如果 $n$ 个元素\textbf{互不相同}，分组数是 $k^n$（每个元素独立选一组）或 $\dbinom{n+k-1}{k-1}$ 乘以组内排列，与插板法完全不同。读题时先看“$n$ 个相同的球”还是“$n$ 个不同的人”。
\end{pitfall}

\section{圆排列与多重集}
\subsection{圆排列}
$n$ 个人围成一圈，排列数为
\[
Q_n=\frac{A_n^n}{n}=(n-1)!,
\]
因为一个圆圈从 $n$ 个不同位置断开会得到 $n$ 个不同的队列。部分圆排列（选 $r$ 个围成一圈）为
\[
Q_n^r=\frac{A_n^r}{r}=\frac{n!}{r\,(n-r)!}.
\]
若再要求“翻转后相同”（如项链），还要除以 $2$。

\subsection{多重集的排列数}
设多重集有 $n_1$ 个 $a_1$、$n_2$ 个 $a_2$、……、$n_k$ 个 $a_k$，总数 $n=\sum n_i$。全排列数为
\[
\binom{n}{n_1,n_2,\ldots,n_k}=\frac{n!}{n_1!\,n_2!\cdots n_k!},
\]
这叫\textbf{多项式系数}。道理很简单：先当作 $n$ 个不同元素全排列，再把相同元素之间的 $n_i!$ 种排列视为同一种。

\subsection{多重集的组合数}
从上述多重集中选 $r$ 个元素（$r<n_i$ 对一切 $i$ 成立，即每种都取不完），方案数就是 $x_1+\cdots+x_k=r$ 的非负整数解数：
\[
\binom{r+k-1}{k-1}.
\]
若每种有上限 $x_i\le n_i$，则要用第 3 章的容斥原理。

\begin{tip}[三个“多重”不要混]
\textbf{多重集的排列数}（$\frac{n!}{\prod n_i!}$，计顺序）、\textbf{多重集的组合数}（$\binom{r+k-1}{k-1}$，不计顺序）、\textbf{多项式系数}（$\binom{n}{n_1,\ldots,n_k}$，就是多重集排列数的另一种记号）。三者中前两个是最容易混的。
\end{tip}

\section{分层例题与解析}
\begin{problem}{例 1：乘法原理}
一个密码由 $3$ 个大写字母后跟 $3$ 个数字组成。若字母不许重复、数字可以重复，有多少个密码？
\end{problem}
\hint 分两步：字母部分与数字部分，各自用乘法原理，再相乘。

\sol 字母部分 $26\cdot25\cdot24$，数字部分 $10^3$，共
\[
26\cdot25\cdot24\cdot10^3=15\,600\,000.
\]

\begin{problem}{例 2：组合与排列的界线}
从 $10$ 名选手中选 $3$ 人组成代表队，再指定一名队长，有多少种方案？若改为“选 $3$ 人并排成一列拍照”呢？
\end{problem}
\hint 第一问先组合再选队长；第二问直接排列。

\sol 第一问 $\dbinom{10}{3}\cdot3=120\cdot3=360$。第二问 $A_{10}^3=720$。两问不同：第一问中另外两人没有身份区别，第二问中三人位置不同。\textbf{先问“顺序算不算”}，答案自然分开。

\begin{problem}{例 3：插板法}
方程 $x_1+x_2+x_3+x_4=20$ 有多少组正整数解？多少组非负整数解？
\end{problem}
\hint 正整数解每组至少 $1$；非负解借 $4$ 个元素。

\sol 正整数解 $\dbinom{20-1}{4-1}=\dbinom{19}{3}=969$。非负整数解 $\dbinom{20+4-1}{4-1}=\dbinom{23}{3}=1771$。

\begin{problem}{例 4：不相邻的选择}
从 $1,2,\ldots,10$ 中选 $3$ 个数，任意两个不相邻，有多少种选法？
\end{problem}
\hint 设选出的数为 $a<b<c$，令 $b'=b-1$、$c'=c-2$ 去掉间隔。

\sol 设 $a<b<c$ 且 $b-a\ge2$、$c-b\ge2$。令 $a'=a$、$b'=b-1$、$c'=c-2$，则 $1\le a'<b'<c'\le8$，这是一个从 $8$ 个数里选 $3$ 的组合，共 $\dbinom83=56$ 种。一般地，从 $n$ 个数里选 $k$ 个互不相邻的数有 $\dbinom{n-k+1}{k}$ 种。

\begin{problem}{统一练习：分配问题}
把 $12$ 本相同的书分给 $4$ 个同学，每人至少 $1$ 本，有多少种分法？若这 $12$ 本书互不相同呢？
\end{problem}
\hint 第一问用插板法；第二问每本书独立选人，再用容斥减去“有人没书”的情形。

\sol 第一问：$x_1+\cdots+x_4=12$ 的正整数解数为
\[
\binom{12-1}{4-1}=\binom{11}{3}=165.
\]

第二问：每本书有 $4$ 种选择，共 $4^{12}$ 种；减去“某个人一本书都没分到”的情形。设 $A_i$ 为“第 $i$ 个人没书”，则
\[
\left|\overline{A_1}\cap\overline{A_2}\cap\overline{A_3}\cap\overline{A_4}\right|
=\sum_{j=0}^{4}(-1)^j\binom4j(4-j)^{12}.
\]
其中 $j=4$ 一项为 $0^{12}=0$（$12$ 本书不可能谁都不给），于是
\[
4^{12}-4\cdot3^{12}+6\cdot2^{12}-4\cdot1^{12}
=16\,777\,216-2\,125\,764+24\,576-4=14\,676\,024.
\]

\begin{teach}[从“至少 $1$ 本”到“至少 $2$ 本”]
若把条件改成“每人至少 $2$ 本”：相同书的情形令 $x_i'=x_i-2$，化为 $\sum x_i'=4$，答案 $\dbinom73=35$；但\textbf{不同}书的情形再也不能用一次容斥解决，需要用生成函数或按“先给每人 $2$ 本再任意分剩下 $4$ 本”的组合分解。后者留到第五部分系统处理。
\end{teach}

\chapter{二项式定理与组合恒等式}
\begin{chapterintro}[章节导读]
$(1+x)^n$ 是组合数学里最重要的一个多项式：它的第 $k$ 项系数就是 $\binom nk$。把“选法条数”看成“多项式系数”以后，乘法变成了卷积、求导变成了加权求和、代入特殊值变成了一整类恒等式。本章建立这套语言，它是第五部分生成函数的直接前身。
\end{chapterintro}

\section{二项式定理}
\begin{theorem}[二项式定理]
对任意非负整数 $n$ 与实数 $x,y$，
\[
(x+y)^n=\sum_{k=0}^{n}\binom nk x^{n-k}y^{k}.
\]
\end{theorem}
\begin{proof}[组合证明]
$(x+y)^n$ 是 $n$ 个 $(x+y)$ 相乘。展开后的每一项要从每个因子里选 $x$ 或 $y$。若恰有 $k$ 个因子选了 $y$（其余 $n-k$ 个选 $x$），得到 $x^{n-k}y^k$；选哪 $k$ 个因子有 $\dbinom nk$ 种，故系数为 $\dbinom nk$。
\end{proof}

\begin{theorem}[多项式定理]
设 $n$ 为非负整数，则
\[
(x_1+x_2+\cdots+x_t)^n
=\sum_{\substack{n_1,\ldots,n_t\ge0\\ n_1+\cdots+n_t=n}}
\binom{n}{n_1,n_2,\ldots,n_t}x_1^{n_1}x_2^{n_2}\cdots x_t^{n_t},
\]
其中 $\dbinom{n}{n_1,\ldots,n_t}=\dfrac{n!}{n_1!\cdots n_t!}$ 是多项式系数。并且
\[
\sum_{\substack{n_1+\cdots+n_t=n}}\binom{n}{n_1,\ldots,n_t}=t^n.
\]
\end{theorem}
\begin{proof}[组合证明]
展开后的每一项对应一个“函数”$f:\{1,\ldots,n\}\to\{1,\ldots,t\}$，$f(i)$ 记录第 $i$ 个因子里选了哪个变量。若 $n_j=|f^{-1}(j)|$，则该项为 $x_1^{n_1}\cdots x_t^{n_t}$，而有同样计数向量的函数有 $\dbinom{n}{n_1,\ldots,n_t}$ 个（先排 $n$ 个位置，再按组除）。函数总数为 $t^n$，即得求和式。
\end{proof}

\begin{tip}[生成函数视角]
把 $\sum_k\binom nk x^k$ 记作 $(1+x)^n$ 的\textbf{系数列}。这个视角的用处：两个多项式相乘时，系数做卷积
\[
\left(\sum_k a_kx^k\right)\left(\sum_k b_kx^k\right)=\sum_k\left(\sum_{i+j=k}a_ib_j\right)x^k.
\]
于是“$(1+x)^m(1+x)^n=(1+x)^{m+n}$”翻译过来就是范德蒙恒等式。第五部分会把这一套完整展开。
\end{tip}

\section{常用组合恒等式}
\begin{theorem}[范德蒙恒等式]
对非负整数 $m,n,k$，
\[
\sum_{i=0}^{k}\binom ni\binom m{k-i}=\binom{m+n}{k}.
\]
特别地，取 $m=n=k$ 得
\[
\sum_{i=0}^{n}\binom ni^2=\binom{2n}{n}.
\]
\end{theorem}
\begin{proof}[生成函数证明]
$(1+x)^m(1+x)^n=(1+x)^{m+n}$。左边 $x^k$ 的系数是 $\sum_{i+j=k}\binom mi\binom nj=\sum_{i=0}^{k}\binom ni\binom m{k-i}$，右边是 $\dbinom{m+n}{k}$。取 $m=n=k$ 并注意 $\binom ni=\binom n{n-i}$，第二个等式是直接推论。
\end{proof}

\begin{theorem}[朱世杰恒等式（冰球棍）]
\[
\sum_{l=0}^{n}\binom lk=\binom{n+1}{k+1}\qquad(k\ge0).
\]
\end{theorem}
\begin{proof}[组合证明]
右边数的是“从 $\{a_1,\ldots,a_{n+1}\}$ 中选 $k+1$ 个元素”的子集数。按\textbf{最大元素}分类：若最大元素是 $a_{l+1}$，则其余 $k$ 个元素要从 $\{a_1,\ldots,a_l\}$ 中选，有 $\dbinom lk$ 种。对 $l=k,\ldots,n$ 求和即得左边。
\end{proof}

\begin{theorem}[带权和的两条恒等式]
\[
\sum_{i=0}^{n}i\binom ni=n2^{n-1},\qquad
\sum_{i=0}^{n}i^2\binom ni=n(n+1)2^{n-2}.
\]
\end{theorem}
\begin{proof}[求导证明]
由二项式定理 $\sum_i\binom nix^i=(1+x)^n$。两边对 $x$ 求导：
\[
\sum_i i\binom nix^{i-1}=n(1+x)^{n-1}.
\]
令 $x=1$ 得第一条。再求导一次并令 $x=1$：
\[
\sum_i i(i-1)\binom ni=n(n-1)2^{n-2},
\]
而 $\sum i^2\binom ni=\sum i(i-1)\binom ni+\sum i\binom ni=n(n-1)2^{n-2}+n2^{n-1}=n(n+1)2^{n-2}$。
\end{proof}

\begin{teach}[“按最大元素分类”是一个通用套路]
朱世杰恒等式的证明按最大元素分类，这个思路可以推广到几乎所有“上标求和”的恒等式。让学生把 $\sum_{l}\binom lk$ 讲成“数 $k+1$ 元子集，看谁最大”，比记住公式有用得多。
\end{teach}

\section{二项式反演}
\begin{theorem}[二项式反演]
设 $f_n$ 表示“恰好使用 $n$ 个不同元素形成某结构的方案数”，$g_n$ 表示“从 $n$ 个不同元素中任选若干个（可以 $0$ 个）形成该结构的总方案数”，则
\[
g_n=\sum_{i=0}^{n}\binom nif_i,
\qquad
f_n=\sum_{i=0}^{n}(-1)^{n-i}\binom nig_i.
\]
\end{theorem}
\begin{proof}[代入验证]
把 $g_i=\sum_{j}\binom ijf_j$ 代入右边：
\[
\sum_{i=0}^{n}(-1)^{n-i}\binom ni\sum_{j=0}^{i}\binom ijf_j
=\sum_{j=0}^{n}f_j\sum_{i=j}^{n}(-1)^{n-i}\binom ni\binom ij.
\]
用恒等式 $\dbinom ni\dbinom ij=\dbinom nj\dbinom{n-j}{i-j}$，内层和变为
\[
\binom nj\sum_{i=j}^{n}(-1)^{n-i}\binom{n-j}{i-j}
=\binom nj\sum_{t=0}^{n-j}(-1)^{n-j-t}\binom{n-j}{t}.
\]
由 $\sum_t(-1)^t\binom{N}{t}=\iva{N=0}$，内层在 $n=j$ 时为 $1$、在 $n>j$ 时为 $0$。故整式和为 $f_n$。
\end{proof}

\begin{tip}[二项式反演与莫比乌斯反演的关系]
第二部分第 8 章的\textbf{莫比乌斯反演}处理的是“按约数”的偏序，二项式反演处理的是“按子集”的偏序。两者形式相同、场景不同：看到“$g_n=\sum\binom nif_i$”用二项式反演，看到“$g(n)=\sum_{d\mid n}f(d)$”用莫比乌斯反演。第五部分会用生成函数把两者统一处理。
\end{tip}

\section{分层例题与解析}
\begin{problem}{例 1：用二项式定理求系数}
求 $(2x-1)^6$ 中 $x^3$ 的系数。
\end{problem}
\hint 直接套二项式定理，$x=2x$、$y=-1$。

\sol 通项为 $\dbinom6k(2x)^{6-k}(-1)^k$。要 $x^3$ 需 $6-k=3$，即 $k=3$，系数为
\[
\binom63\cdot2^3\cdot(-1)^3=20\cdot8\cdot(-1)=-160.
\]

\begin{problem}{例 2：范德蒙恒等式}
计算 $\sum_{i=0}^{5}\binom{5}{i}\binom{5}{5-i}$。
\end{problem}
\hint 这就是 $\sum_i\binom5i\binom5{5-i}$，令 $m=n=k=5$。

\sol 由范德蒙恒等式的特殊形式，$\sum_i\binom5i\binom5{5-i}=\sum_i\binom5i^2=\binom{10}{5}=252$。

\begin{problem}{例 3：组合数什么时候为零}
计算 $\dbinom30+\dbinom31+\dbinom32+\cdots+\dbinom{3}{10}$。
\end{problem}
\hint 先看每一项是否为 $0$；$\binom nk=0$ 当 $k>n$。

\sol $\binom nk=0$（$k>n$），故只有前四项非零：
\[
\binom30+\binom31+\binom32+\binom33=1+3+3+1=8.
\]
这与 $(1+1)^3=2^3=8$ 一致。这个例子提醒：\textbf{求和上下界要先与组合数的有效范围对齐}，否则容易把空项当成非零项反复计算。

\begin{problem}{例 4：二项式反演}
$n$ 个元素的集合有多少个真子集（不含自身）？用二项式反演验证。
\end{problem}
\hint 令 $f_n$ 为“恰好用 $n$ 个元素”的方案数（只有 $1$ 种：全选），$g_n$ 为总方案数。

\sol $f_n=1$ 对一切 $n$，故
\[
g_n=\sum_{i=0}^{n}\binom ni\cdot1=2^n,
\]
正是子集总数。反过来
\[
f_n=\sum_{i=0}^{n}(-1)^{n-i}\binom ni2^i=(2-1)^n=1,
\]
自洽。这个练习用来说明：二项式反演的两条公式互为逆运算，验证时先代入一条再代回另一条。

\chapter{容斥原理}
\begin{chapterintro}[章节导读]
“至少满足一个条件”的对象有多少个？直接数会重复，容斥原理给出系统的加减法。它是组合数学里最强大的单一工具：从错位排列到欧拉函数，从格路计数到概率里的并集，全都落在它身上。第二部分第 8 章的莫比乌斯反演，本质上就是容斥原理的约数版本。
\end{chapterintro}

\section{容斥原理的形式}
\begin{theorem}[容斥原理]
设 $A_1,\ldots,A_n$ 是有限集，则
\[
\left|\bigcup_{i=1}^{n}A_i\right|
=\sum_{\varnothing\ne S\subseteq[n]}(-1)^{|S|+1}\left|\bigcap_{i\in S}A_i\right|.
\]
等价地，用补集写：
\[
\left|\overline{A_1}\cap\cdots\cap\overline{A_n}\right|
=\sum_{S\subseteq[n]}(-1)^{|S|}\left|\bigcap_{i\in S}A_i\right|.
\]
\end{theorem}
\begin{proof}[按元素被数的次数验证]
固定一个元素 $x$，设它恰好属于 $A_1,\ldots,A_n$ 中的 $r$ 个（$r\ge1$）。左边只把它数 $1$ 次。右边：含 $x$ 的交集 $\bigcap_{i\in S}A_i$ 恰对应 $S\subseteq T$（$T$ 是那 $r$ 个指标组成的集合），共 $2^r$ 项，其贡献为
\[
\sum_{S\subseteq T}(-1)^{|S|+1}=-\sum_{j=0}^{r}\binom rj(-1)^j=-\iva{r=0}=1\qquad(r\ge1).
\]
恰好 $1$ 次。若 $r=0$（$x$ 不属于任何 $A_i$），左边为 $0$，右边所有交集都不含 $x$，贡献也是 $0$。两边对每个元素的计数相同，故相等。
\end{proof}

\begin{teach}[这个证明值得反复讲]
容斥原理的证明不化简任何代数式，只做一件事：\textbf{看一个固定元素被数了几次}。这个手法在第二部分第 5 章证明 LTE 的“唯一最低赋值项”、第三部分第 5 章的抽屉原理里都会再现。建议让学生在黑板上把 $n=3$ 的情形逐个数一遍，再回到一般证明。
\end{teach}

\section{错位排列}
\begin{definition}[错位排列]
$n$ 个元素的一个排列，若每个元素都不在原来的位置上，称为\textbf{错位排列}。错位排列数记作 $D_n$。
\end{definition}
\begin{theorem}[错位排列数]
\[
D_n=n!\sum_{k=0}^{n}\frac{(-1)^k}{k!}.
\]
特别地，$D_0=1$、$D_1=0$、$D_2=1$、$D_3=2$、$D_4=9$、$D_5=44$。且
\[
\lim_{n\to\infty}\frac{D_n}{n!}=\frac1{\mathrm e}.
\]
\end{theorem}
\begin{proof}
设 $A_i$ 为“第 $i$ 个元素留在原位”的排列集合。则 $|A_i|=(n-1)!$，更一般地 $|A_{i_1}\cap\cdots\cap A_{i_k}|=(n-k)!$。由容斥，
\[
D_n=n!-\left|\bigcup_i A_i\right|
=n!-\sum_{k=1}^{n}(-1)^{k+1}\binom nk(n-k)!
=n!\sum_{k=0}^{n}\frac{(-1)^k}{k!}.
\]
极限由 $\mathrm e^{-1}=\sum_k(-1)^k/k!$ 得到。
\end{proof}

\begin{tip}[错位排列的递推]
$D_n$ 也满足递推 $D_n=(n-1)(D_{n-1}+D_{n-2})$（$n\ge2$）。组合解释：第 $1$ 个元素放到位置 $j$（$n-1$ 种选择），再分“第 $j$ 个元素回到位置 $1$”与“不回去”两种，分别对应 $D_{n-2}$ 与 $D_{n-1}$。两种推导都要会。
\end{tip}

\section{容斥的三种典型应用}
\subsection{有上限的多重集组合}
从 $n_1$ 个 $a_1$、……、$n_k$ 个 $a_k$ 中选 $r$ 个，要求 $x_i\le n_i$。全集是非负整数解 $\dbinom{r+k-1}{k-1}$；“违反第 $i$ 个上限”即 $x_i\ge n_i+1$，令 $x_i'=x_i-n_i-1$ 后化为非负整数解，可逐项计算。

\subsection{欧拉函数的容斥表达式}
第二部分第 2 章会证明
\[
\varphi(n)=n\prod_{p\mid n}\left(1-\frac1p\right),
\]
其推导正是容斥：$\varphi(n)$ 数是 $1,\ldots,n$ 中与 $n$ 互质的整数个数，即“不被任何素因子整除”的个数。

\subsection{格路与障碍物}
第三部分第 6 章的反射原理，是容斥原理在“路径被挡住”场景下的精确形式：把坏路径双射到另一类好数的路径，从而减去。

\section{分层例题与解析}
\begin{problem}{例 1：基本容斥}
某班 $50$ 人，参加数学小组的 $30$ 人，参加编程小组的 $25$ 人，两个都参加的 $15$ 人。两个都没参加的有多少人？
\end{problem}
\hint 用 $|A\cup B|=|A|+|B|-|A\cap B|$。

\sol $|A\cup B|=30+25-15=40$，故都没参加的有 $50-40=10$ 人。

\begin{problem}{例 2：错位排列}
$5$ 个人各写一张纸条放入帽中，再随机取出一张（不是自己的）。有多少种取法？
\end{problem}
\hint 这就是 $D_5$。

\sol $D_5=5!\left(1-1+\dfrac12-\dfrac16+\dfrac1{24}-\dfrac1{120}\right)=120\cdot\dfrac{44}{120}=44$。

\begin{problem}{例 3：容斥计数}
求 $1$ 到 $1000$ 中能被 $6$、$10$ 或 $15$ 至少一个整除的整数个数。
\end{problem}
\hint 用 $|A\cup B\cup C|=|A|+|B|+|C|-|A\cap B|-|A\cap C|-|B\cap C|+|A\cap B\cap C|$；每一处交集都先算最小公倍数。

\sol 记 $A,B,C$ 分别为被 $6,10,15$ 整除的数集。
\[
|A|=\floor{\frac{1000}{6}}=166,\quad |B|=\floor{\frac{1000}{10}}=100,\quad |C|=\floor{\frac{1000}{15}}=66.
\]
交集用最小公倍数：
\[
\lcm(6,10)=30,\quad \lcm(6,15)=30,\quad \lcm(10,15)=30,\quad \lcm(6,10,15)=30,
\]
故三个两两交集各为 $\floor{1000/30}=33$，三者交集也是 $33$。于是
\[
|A\cup B\cup C|=166+100+66-33-33-33+33=266.
\]

\begin{pitfall}[交集用的是最小公倍数]
“能被 $a$ 且能被 $b$ 整除”等价于“能被 $\lcm(a,b)$ 整除”。本题三对最小公倍数恰好都是 $30$，但若换成 $6$ 与 $10$、$14$ 与 $21$ 之类的组合，$\lcm$ 各不相同，\textbf{每一处交集都要单独算}。
\end{pitfall}

\begin{problem}{例 4：容斥求欧拉函数}
用容斥原理求 $\varphi(30)$，并与公式 $n\prod_{p\mid n}(1-\frac1p)$ 对照。
\end{problem}
\hint $30=2\cdot3\cdot5$；在 $1,\ldots,30$ 中数不被 $2,3,5$ 整除的数。

\sol 在 $1,\ldots,30$ 中，被 $2$ 整除的 $15$ 个、被 $3$ 整除的 $10$ 个、被 $5$ 整除的 $6$ 个；被 $2$ 与 $3$ 整除的 $5$ 个、被 $2$ 与 $5$ 整除的 $3$ 个、被 $3$ 与 $5$ 整除的 $2$ 个；被 $2,3,5$ 同时整除的 $1$ 个。于是
\[
\varphi(30)=30-(15+10+6)+(5+3+2)-1=30-31+10-1=8.
\]
公式给出
\[
30\left(1-\frac12\right)\left(1-\frac13\right)\left(1-\frac15\right)
=30\cdot\frac12\cdot\frac23\cdot\frac45=8,
\]
两者一致。

\begin{teach}[容斥的算术要逐项写]
本题只有 $7$ 个非零项，但每一项的分子分母都容易抄错。讲评时要求学生把 $30/2=15$、$30/3=10$、$30/5=6$、$30/6=5$、$30/10=3$、$30/15=2$、$30/30=1$ 七个数值先列成一列，再代入公式。\textbf{容斥的失分几乎都在算术，不在思路。}
\end{teach}

\chapter{递推关系与经典计数数列}
\begin{chapterintro}[章节导读]
很多组合问题的答案不能直接套公式，但满足一个递推关系：大问题的解由小问题的解拼成。本章处理三条最经典的递推——斐波那契、卡特兰、斯特林——并说明“递推”与“闭式”的关系：能写出递推往往就够了，闭式只是为了更快地算。
\end{chapterintro}

\section{怎么建立递推关系}
建立递推的标准动作是\textbf{按最后一个动作分类}：
\begin{enumerate}
\item 指出“一个合法对象”的最后一步有几种可能；
\item 对每种可能，说明去掉最后一步以后得到的是一个（或几个）同类的小对象；
\item 把各类计数相加或相乘，得到递推式；
\item 写出边界条件。
\end{enumerate}

\begin{pitfall}[重复计数与漏计]
递推式最常见的错误是同一个对象被两条路径数到（重复）或某类对象没有路径能到达（漏计）。写完递推式后，用 $n=1,2,3$ 的小值手工枚举核对，是成本最低的验证。
\end{pitfall}

\section{斐波那契数列}
\begin{definition}[斐波那契数列]
$F_1=F_2=1$，$F_{n+2}=F_{n+1}+F_n$。前几项：
\[
1,1,2,3,5,8,13,21,34,55,89,\ldots
\]
\end{definition}
\begin{theorem}[斐波那契的组合解释]
$F_{n+1}$ 等于用 $1\times1$ 与 $1\times2$ 的方块铺满 $1\times n$ 的长条的方法数。
\end{theorem}
\begin{proof}
设铺法数为 $f_n$，$f_0=1$（空条有一种铺法）。按最左边一块分类：若最左是 $1\times1$，剩下 $1\times(n-1)$ 有 $f_{n-1}$ 种；若最左是 $1\times2$，剩下有 $f_{n-2}$ 种。故 $f_n=f_{n-1}+f_{n-2}$，$f_0=1$、$f_1=1$，即 $f_n=F_{n+1}$。
\end{proof}

\begin{tip}[斐波那契的增长速度]
由特征方程 $x^2=x+1$ 得 $F_n=\dfrac{1}{\sqrt5}\left[\left(\dfrac{1+\sqrt5}{2}\right)^n-\left(\dfrac{1-\sqrt5}{2}\right)^n\right]$。于是 $F_n=\Theta(\varphi^n)$，其中 $\varphi=\dfrac{1+\sqrt5}{2}\approx1.618$ 是黄金比。由此可得实用估计：$F_{45}\approx1.1\times10^9$、$F_{90}\approx2.9\times10^{18}$，已经贴近 \texttt{long long} 上限。
\end{tip}

\section{卡特兰数}
\begin{definition}[卡特兰数]
\[
C_n=\frac1{n+1}\binom{2n}{n}=\binom{2n}{n}-\binom{2n}{n+1}.
\]
前几项：$C_0=1,C_1=1,C_2=2,C_3=5,C_4=14,C_5=42,C_6=132$。
\end{definition}
\begin{theorem}[卡特兰数的递推]
\[
C_0=1,\qquad C_{n+1}=\sum_{i=0}^{n}C_iC_{n-i}\quad(n\ge0).
\]
\end{theorem}
\begin{proof}[以括号序列为例]
$C_n$ 等于 $n$ 对括号的正确匹配方案数。设 $f_n$ 为该数，$f_0=1$。按\textbf{第一个左括号与谁匹配}分类：若它与第 $2i+2$ 个字符处的右括号匹配（$0\le i\le n-1$），则这对括号内部有 $i$ 对括号（$f_i$ 种），外部有 $n-1-i$ 对（$f_{n-1-i}$ 种），共 $f_if_{n-1-i}$ 种。求和得
\[
f_n=\sum_{i=0}^{n-1}f_if_{n-1-i},
\]
即 $f_{n+1}=\sum_{i=0}^{n}f_if_{n-i}$。再用生成函数或归纳可证 $f_n=C_n$。
\end{proof}

\begin{tip}[卡特兰数的家族]
以下问题的答案都是 $C_n$：$n$ 对括号的正确匹配；$n$ 个结点能形成的不同二叉树的个数；从 $(0,0)$ 到 $(n,n)$ 不越过对角线的单调格路数；$n$ 个元素依次入栈、全部出栈的合法顺序数；“$n+2$ 边形用不相交对角线分成三角形”的方法数。第三部分第 6 章会详细处理格路版本。
\end{tip}

\section{斯特林数}
\subsection{第二类斯特林数}
\begin{definition}[第二类斯特林数]
$S(n,k)$（也记作 $\begin{Bmatrix}n\\k\end{Bmatrix}$）表示把 $n$ 个\textbf{不同}的元素分成 $k$ 个\textbf{非空无标号}子集的方法数。
\end{definition}
\begin{theorem}[第二类斯特林数的递推]
\[
S(n,k)=k\,S(n-1,k)+S(n-1,k-1),
\]
边界 $S(0,0)=1$、$S(n,0)=0$（$n\ge1$）、$S(0,k)=0$（$k\ge1$）。
\end{theorem}
\begin{proof}[按第 $n$ 个元素分类]
把 $n$ 个元素分成 $k$ 个非空无标号子集。看第 $n$ 个元素：
\begin{itemize}
\item 它单独成组：其余 $n-1$ 个元素分成 $k-1$ 组，共 $S(n-1,k-1)$ 种；
\item 它加入已有的某一组：先把前 $n-1$ 个元素分成 $k$ 组（$S(n-1,k)$ 种），再把第 $n$ 个元素放入 $k$ 个组中的一个（$k$ 种选择），共 $kS(n-1,k)$ 种。
\end{itemize}
两类互斥，相加即得。
\end{proof}

\begin{tip}[“无标号”与“有标号”只差一个 $k!$]
若 $k$ 个盒子\textbf{有标号}，方案数是 $k!\,S(n,k)$。这是斯特林数应用中最常见的陷阱：读题时先看“盒子有没有编号”。
\end{tip}

\subsection{第一类斯特林数}
\begin{definition}[第一类斯特林数]
$s(n,k)$（带符号）与 $c(n)=\begin{bmatrix}n\\k\end{bmatrix}$（不带符号）表示把 $n$ 个不同元素排成 $k$ 个不相交的\textbf{圆排列}的方法数。
\end{definition}
\begin{theorem}[第一类斯特林数的递推]
\[
\begin{bmatrix}n\\k\end{bmatrix}=(n-1)\begin{bmatrix}n-1\\k\end{bmatrix}+\begin{bmatrix}n-1\\k-1\end{bmatrix},
\]
边界 $\begin{bmatrix}0\\0\end{bmatrix}=1$。它们还是下降阶乘的系数：
\[
x(x-1)(x-2)\cdots(x-n+1)=\sum_{k=0}^{n}s(n,k)x^k.
\]
\end{theorem}
\begin{proof}[递推的组合解释]
看第 $n$ 个元素：若它单独成一个圆排列，其余 $n-1$ 个排成 $k-1$ 个圆，有 $\begin{bmatrix}n-1\\k-1\end{bmatrix}$ 种；若它插入已有的某个圆排列，该圆有 $m$ 个元素时有 $m$ 个插入位置，对所有圆求和共 $n-1$ 个位置，故有 $(n-1)\begin{bmatrix}n-1\\k\end{bmatrix}$ 种。
\end{proof}

\section{分层例题与解析}
\begin{problem}{例 1：铺砖}
用 $1\times1$ 与 $1\times2$ 的方块铺满 $1\times6$ 的长条，有多少种方法？
\end{problem}
\hint 就是 $F_7$。

\sol $F_7=13$。可以直接列：$f_0=1,f_1=1,f_2=2,f_3=3,f_4=5,f_5=8,f_6=13$。

\begin{problem}{例 2：卡特兰}
$3$ 对括号有多少种正确匹配？$4$ 个结点能形成多少种不同的二叉树？
\end{problem}
\hint 都是 $C_3$ 与 $C_4$。

\sol $C_3=\dfrac16\dbinom63=\dfrac{20}{6}=5$；$C_4=\dfrac15\dbinom84=\dfrac{70}{5}=14$。

枚举验证 $3$ 对括号的 $5$ 种匹配：
\[
\texttt{((()))},\quad \texttt{(()())},\quad \texttt{(())()},\quad \texttt{()(())},\quad \texttt{()()()}.
\]
只有这 $5$ 种，不多不少。

\begin{teach}[枚举必须做到“不重不漏”]
让学生自己列出 $3$ 对括号的全部匹配，几乎一定会出现重复或遗漏。正确的做法是按第 4 章的分类：第一个左括号与谁配对。若它与第 $2i+2$ 个字符处的右括号配对（$i=0,1,2$），则内部有 $i$ 对、外部有 $2-i$ 对，于是
\[
f_2=f_0f_1+f_1f_0+f_2f_0=3,\qquad
f_3=f_0f_2+f_1f_1+f_2f_0=2+1+2=5.
\]
\textbf{按分类枚举才不会重、不会漏。}上面列出的 $5$ 种正是按“最外层括号内含 $0$ 对、$1$ 对、$2$ 对”分三类的代表。
\end{teach}

\begin{problem}{例 3：第二类斯特林数}
把 $5$ 个人分成 $3$ 个非空小组（组无编号），有多少种分法？
\end{problem}
\hint 就是 $S(5,3)$；用递推算。

\sol 列表计算：
\[
S(1,1)=1;\quad S(2,1)=1,\ S(2,2)=1;
\]
\[
S(3,1)=1,\ S(3,2)=2\cdot1+1=3,\ S(3,3)=1;
\]
\[
S(4,1)=1,\ S(4,2)=2\cdot3+1=7,\ S(4,3)=3\cdot1+3=6,\ S(4,4)=1;
\]
\[
S(5,1)=1,\ S(5,2)=2\cdot7+1=15,\ S(5,3)=3\cdot6+7=25.
\]
答案 $25$。若小组有编号，则 $3!\cdot25=150$。

\begin{problem}{例 4：递推的建立}
一条楼梯有 $n$ 级，每次可以上一级或两级，有多少种上楼方法？若每次可以上一、二或三级呢？
\end{problem}
\hint 按最后一步分类。

\sol 设 $f_n$ 为 $n$ 级楼梯的上法数，$f_0=1$。每次一或两级：$f_n=f_{n-1}+f_{n-2}$，即 $f_n=F_{n+1}$。每次一、二或三级：$f_n=f_{n-1}+f_{n-2}+f_{n-3}$，$f_0=1,f_1=1,f_2=2$，于是 $f_3=4,f_4=7,f_5=13$，这就是\textbf{ Tribonacci 数列}。

\begin{problem}{统一练习：卡特兰的又一副面孔}
$n$ 个元素依次入栈，求所有可能的出栈序列数。用 $n=3$ 验证答案为 $C_3$。
\end{problem}
\hint 把入栈记作左括号、出栈记作右括号；出栈数不能超过入栈数。

\sol 入栈序列记作 $1$，出栈记作 $-1$。任意前缀中 $1$ 的个数不少于 $-1$ 的个数，这正是括号匹配的条件，故答案为 $C_n$。$n=3$ 时有 $5$ 种：$123,132,213,231,321$。注意 $312$ 不可能（$3$ 第一个出栈要求 $1,2$ 已在栈中，那么出栈顺序只能是 $321$ 或 $23\ldots$，读者可自行核对）。

\begin{teach}[递推章的教学要点]
\begin{enumerate}
\item 每道题都要求学生先写“按最后一个动作分类”的那句话，再写递推式。
\item 递推式写完后必须用 $n=1,2,3$ 手工核对；对不上就重写分类。
\item 闭式公式（比奈公式、$\frac1{n+1}\binom{2n}{n}$）只作为“算得快”的工具，不要求会推导。
\end{enumerate}
\end{teach}

\chapter{抽屉原理与极端原理}
\begin{chapterintro}[章节导读]
前四章都在“数”，本章开始“证存在”。抽屉原理（鸽巢原理）说：东西比格子多，就有一格不止一个。它是组合数学里最朴素也最强大的工具，而它的两个常见变体——平均值原理与极端原理——分别处理“平均”和“最值”型的存在性论证。
\end{chapterintro}

\section{抽屉原理}
\begin{theorem}[抽屉原理]
把 $n+1$ 个物体放进 $n$ 个抽屉，至少有一个抽屉里有不少于 $2$ 个物体。
\end{theorem}
\begin{proof}[反证]
若每个抽屉至多 $1$ 个物体，则总数至多 $n$，与 $n+1$ 矛盾。
\end{proof}

\begin{theorem}[抽屉原理的加强形式]
把 $mn+1$ 个物体放进 $n$ 个抽屉，至少有一个抽屉里有不少于 $m+1$ 个物体。等价地：若 $n$ 个整数的平均值大于 $a$，则至少有一个整数大于 $a$。
\end{theorem}

\begin{teach}[抽屉原理的三步讲法]
\begin{enumerate}
\item \textbf{造抽屉}：把对象按某个性质分类，说明类别数少于对象数。这是最难的一步——题目从不告诉你抽屉是什么。
\item \textbf{数物体}：说明物体数比抽屉数多（或用加强形式给出 $m$）。
\item \textbf{得结论}：把“某类里有两个物体”翻译回原问题的语言。
\end{enumerate}
第一步没有通法，只能靠题目积累。常见的造抽屉方式有：按余数分类（模 $n$ 的剩余类）、按数值区间分块、按图的度数分层、按坐标奇偶性分类。
\end{teach}

\section{平均值原理}
\begin{theorem}[平均值原理]
若 $n$ 个实数的平均值为 $\bar a$，则至少有一个数不小于 $\bar a$，也至少有一个数不大于 $\bar a$。
\end{theorem}
\begin{proof}[反证]
若所有数都小于 $\bar a$，则平均数小于 $\bar a$，矛盾。大于的情形同理。
\end{proof}
平均值原理是抽屉原理的“加权版”：把 $n$ 个数看成物体、把“是否达到平均”看成两个抽屉。它在图论里最常用：度数平均值 $\dfrac{2m}{n}$ 保证存在度数不小于平均值的顶点。

\section{极端原理}
\begin{theorem}[极端原理]
在一组有限的对象中，取“最大”“最小”“最长”“最短”的那个，利用它的极端性导出矛盾或构造。这与第一部分第 1 章的“最小反例法”是同一件事。
\end{theorem}

\begin{tip}[什么时候用极端原理]
当题目问“证明存在一个……”，而直接构造很困难时，取一个极端对象（例如面积最小的反例、距离最远的一对点、字典序最小的解），它的极端性往往正好提供你需要的那条不等式。第二部分第 15 章的韦达跳跃，就是取“最小的解”然后造出更小的解来完成递降。
\end{tip}

\section{染色与覆盖}
染色是抽屉原理的几何版本：把对象染成少数几种颜色，利用“同色对象必有某种关系”得出结论。经典模式：
\begin{itemize}
\item \textbf{二染色}：把棋盘染成黑白两色，数两种颜色的格子数差。
\item \textbf{模 $k$ 染色}：把整数按模 $k$ 分类，用于处理周期性与整除。
\item \textbf{覆盖论证}：用若干个小图形覆盖大图形，比较面积，得出“至少有两个小图形重叠”。
\end{itemize}

\begin{pitfall}[“至少”与“至多”别写反]
“$n+1$ 个物体放进 $n$ 个抽屉”得到的是“\textbf{至少有一个}抽屉不少于 $2$ 个”，不是“所有抽屉都不少于 $2$ 个”，也不是“至多有一个”。写结论时把量词写全。
\end{pitfall}

\section{分层例题与解析}
\begin{problem}{例 1：基本抽屉}
证明：任意 $13$ 个人中，至少有两个人的生日在同一月份。
\end{problem}
\hint $12$ 个月是抽屉。

\sol $13$ 个人放进 $12$ 个月份，由抽屉原理至少有一个月份有两人以上。注意\textbf{抽屉是月份、物体是人}，不要说反。

\begin{problem}{例 2：加强形式}
证明：任意 $n+1$ 个正整数中，至少有两个数的差能被 $n$ 整除。
\end{problem}
\hint 按模 $n$ 的余数造抽屉。

\sol 两个数的差能被 $n$ 整除 $\iff$ 它们模 $n$ 同余。余数只有 $n$ 种，$n+1$ 个数必有同余的两个。

\begin{problem}{例 3：平均值原理}
证明：任意一群人，认识人数最多的人与最少的人，其认识人数之差不超过总人数减一。（把“认识”视为对称关系。）
\end{problem}
\hint 认识人数是整数，范围在 $0$ 与 $n-1$ 之间。

\sol 设共 $n$ 个人。每个人的认识人数是 $0$ 到 $n-1$ 之间的整数，但 $0$ 与 $n-1$ 不能同时出现（若有人认识所有人，则没有人谁也不认识）。故实际取值至多 $n-1$ 种，极差至多 $n-2$。这个例子用来说明：\textbf{先检查极端值能否同时取到}，是平均值与抽屉问题的常见陷阱。

\begin{problem}{例 4：极端原理}
设 $a_1,\ldots,a_n$ 是正整数，证明其中存在若干个（可以一个）的和能被 $n$ 整除。
\end{problem}
\hint 考虑前缀和 $S_k=a_1+\cdots+a_k$；若有两个前缀和同余则成功，否则 $n$ 个前缀和占据全部 $n$ 个余数类。

\sol 考虑前缀和 $S_0=0,S_1,\ldots,S_n$ 共 $n+1$ 个数。若其中两个模 $n$ 同余，设 $0\le i<j\le n$，则 $S_j-S_i=a_{i+1}+\cdots+a_j$ 被 $n$ 整除。若两两不同余，则 $S_1,\ldots,S_n$ 恰好取遍 $0,\ldots,n-1$ 每个余数一次，特别地某个 $S_k\equiv0\pmod n$，即前 $k$ 项和被 $n$ 整除。两种情况都得结论。

\begin{problem}{例 5：染色}
在 $8\times8$ 的棋盘上任意染黑 $13$ 个格子，证明必有两黑格同行或同列。
\end{problem}
\hint 用两次抽屉，或用行列计数。

\sol 若没有两黑格同行，则每行至多一个黑格，$8$ 行至多 $8$ 个黑格，矛盾。故必有同行。同列的论证完全一样。更强的结论：$13$ 个黑格必有某一行或某一列含不少于 $2$ 个黑格——这正是抽屉原理的直接应用。

\begin{problem}{统一练习：Ramsey 数 $R(3,3)=6$}
证明：任意 $6$ 个人中，必有 $3$ 人两两认识，或 $3$ 人两两不认识。
\end{problem}
\hint 固定一个人 $A$，看 $A$ 与其余 $5$ 人的关系；用抽屉原理得 $A$ 至少认识或不认识 $3$ 人。

\sol 固定 $A$。$A$ 与其余 $5$ 人每人要么认识要么不认识，由抽屉原理，$A$ 至少对其中 $3$ 人有相同的关系。设 $A$ 认识 $B,C,D$（不认识的情形对称）。若 $B,C,D$ 中有两人互相认识，这两人与 $A$ 构成三人两两认识；否则 $B,C,D$ 三人两两不认识。两种情况都成立。

这个证明是 Ramsey 理论的入门例子。它展示了抽屉原理的深层用法：\textbf{先固定一个对象，把“关系”变成“抽屉”}。

\begin{teach}[本章的离堂检查]
让学生独立完成“证明任意 $n+2$ 个正整数中必有两个数的差能被 $n$ 或能被 $n+1$ 整除”这类变体（提示：造 $n+1$ 个抽屉而不是 $n$ 个，或者用加强形式）。做不出就回到例 2 重讲“按余数造抽屉”。
\end{teach}

\chapter{格路计数与反射原理}
\begin{chapterintro}[章节导读]
格路是组合数学里最具体的对象之一：从 $(0,0)$ 出发，每步向右或向上走一格。它的计数看似简单，却能把卡特兰数、选票问题、反射原理串联成一个整体，并直接对应到第二部分数论里的路径分解思想。本章把这条线走完。
\end{chapterintro}

\section{单调格路的基本计数}
\begin{definition}[单调格路]
从 $(0,0)$ 到 $(m,n)$、每步只能向右 $(+1,0)$ 或向上 $(0,+1)$ 的路径叫\textbf{单调格路}。
\end{definition}
\begin{theorem}[单调格路计数]
从 $(0,0)$ 到 $(m,n)$ 的单调格路数为
\[
\binom{m+n}{m}.
\]
\end{theorem}
\begin{proof}
一条路径由 $m+n$ 步组成，其中恰有 $m$ 步向右。选哪 $m$ 步向右有 $\dbinom{m+n}{m}$ 种，其余步自然向上，与路径一一对应。
\end{proof}

\begin{tip}[格路是“序列”的几何化身]
把右步记作 $0$、上步记作 $1$，一条格路就是一个长度为 $m+n$、含 $m$ 个 $0$ 的 $01$ 序列。所以格路计数与“选位置”是同一件事。这个对应在第五部分处理多项式系数时会再次出现。
\end{tip}

\section{不越过对角线的格路与卡特兰数}
\begin{theorem}[不越过对角线的格路]
从 $(0,0)$ 到 $(n,n)$、每步向右或向上、且\textbf{不越过直线 $y=x$}（允许接触）的单调格路数为卡特兰数
\[
C_n=\frac1{n+1}\binom{2n}{n}.
\]
\end{theorem}
\begin{proof}[反射原理]
先数全部路径 $\dbinom{2n}{n}$，再减掉越线的坏路径。一条路径越线，意味着它第一次到达直线 $y=x+1$（即第一次上步比右步多 $1$）。把这一步之后的路径\textbf{关于直线 $y=x+1$ 作镜像}：右步变上步、上步变右步。镜像后路径从 $(0,0)$ 到 $(n-1,n+1)$。

反过来，任意一条从 $(0,0)$ 到 $(n-1,n+1)$ 的路径必然穿过 $y=x+1$，取第一次穿过的位置之后作同样的镜像，就得到一条从 $(0,0)$ 到 $(n,n)$ 的坏路径。于是坏路径与“到 $(n-1,n+1)$ 的路径”一一对应，共 $\dbinom{2n}{n-1}$ 条。故
\[
C_n=\binom{2n}{n}-\binom{2n}{n-1}
=\binom{2n}{n}-\frac n{n+1}\binom{2n}{n}
=\frac1{n+1}\binom{2n}{n}.
\]
\end{proof}

\begin{teach}[反射原理的两步核对]
\begin{enumerate}
\item \textbf{双射}：坏路径与“到 $(n-1,n+1)$ 的路径”之间的映射必须是双射（互为逆）。让学生对 $n=2$ 枚举验证。
\item \textbf{边界}：“不越过”允许接触对角线；“不接触”（严格在下方）是另一个问题，答案为 $C_{n-1}$。两种题面的区别务必讲清。
\end{enumerate}
\end{teach}

\section{选票问题}
\begin{theorem}[选票问题]
一次选举中候选人甲得 $p$ 票、乙得 $q$ 票，$p>q$。若计票过程中甲的票数\textbf{始终严格多于}乙，则计票顺序数为
\[
\frac{p-q}{p+q}\binom{p+q}{p}.
\]
\end{theorem}
\begin{proof}[反射原理]
把每一票记作一步：甲票记右步、乙票记上步。计票顺序就是一条从 $(0,0)$ 到 $(p,q)$ 的单调格路，“甲始终严格多于乙”即路径始终不触碰直线 $y=x$（只能在其下方）。

坏路径是那些触碰 $y=x$ 的。取第一次触碰 $y=x$ 的位置，把之前的路径关于 $y=x$ 镜像，得到一条从 $(0,0)$ 到 $(q,p)$ 的路径；这是一个双射。故好路径数为
\[
\binom{p+q}{p}-\binom{p+q}{q}
=\binom{p+q}{p}-\frac{q}{p}\binom{p+q}{p}
=\frac{p-q}{p+q}\binom{p+q}{p}.
\]
\end{proof}

\begin{tip}[选票问题是卡特兰的一般化]
取 $p=q+1=n+1$、$q=n$，得 $\dfrac{1}{2n+1}\dbinom{2n+1}{n}=C_n$，正是第 4 章的卡特兰数。这个特例说明“不越过”与“始终严格多于”两种边界条件对应不同的题面，用同一个反射原理处理。
\end{tip}

\section{带障碍的格路}
若格路上有若干禁入点，可以用\textbf{容斥原理}：总路径数减去经过禁入点的路径数。若禁入点可以按偏序排列（例如从左下到右上），经过它们的路径要交错的并集，用容斥展开。

\begin{problem}{统一练习：一个禁入点}
从 $(0,0)$ 到 $(4,4)$ 的单调格路中，不经过 $(2,2)$ 的有多少条？
\end{problem}
\hint 总数减去经过 $(2,2)$ 的；经过 $(2,2)$ 的路径 = 到 $(2,2)$ 的路径数 $\times$ 从 $(2,2)$ 到 $(4,4)$ 的路径数。

\sol 总数 $\dbinom84=70$。经过 $(2,2)$ 的路径数为 $\dbinom42\cdot\dbinom42=6\cdot6=36$。故答案为 $70-36=34$。

\section{分层例题与解析}
\begin{problem}{例 1：基本格路计数}
从 $(0,0)$ 到 $(5,3)$ 有多少条单调格路？
\end{problem}
\hint 共 $8$ 步，选 $5$ 步向右。

\sol $\dbinom85=56$。

\begin{problem}{例 2：反射原理}
求从 $(0,0)$ 到 $(4,4)$ 不越过对角线 $y=x$ 的单调格路数。
\end{problem}
\hint 就是 $C_4$。

\sol $C_4=\dfrac15\dbinom84=\dfrac{70}{5}=14$。也可以用反射原理直接算：$\dbinom84-\dbinom8{3}=70-56=14$。

\begin{problem}{例 3：选票问题}
一次选举甲得 $5$ 票、乙得 $3$ 票，求计票过程中甲始终严格多于乙的顺序数。
\end{problem}
\hint 代入公式。

\sol $\dfrac{5-3}{5+3}\dbinom85=\dfrac28\cdot56=14$。可以枚举验证：$5$ 个 $A$、$3$ 个 $B$ 的序列中，每个前缀里 $A$ 的个数都大于 $B$ 的个数，共 $14$ 个。

\begin{problem}{例 4：容斥处理两个禁入点}
从 $(0,0)$ 到 $(5,5)$ 的单调格路，不经过 $(2,2)$ 与 $(3,3)$ 的有多少条？
\end{problem}
\hint 用容斥；经过 $(2,2)$ 的路径数 = 到 $(2,2)$ 的路径数 $\times$ 从 $(2,2)$ 到 $(5,5)$ 的路径数。

\sol 总数 $\dbinom{10}5=252$。记 $A$ 为经过 $(2,2)$ 的路径集，$B$ 为经过 $(3,3)$ 的路径集。
\[
|A|=\binom42\binom{6}{3}=6\cdot20=120,\qquad
|B|=\binom63\binom42=20\cdot6=120.
\]
交集 $A\cap B$：路径同时经过 $(2,2)$ 与 $(3,3)$，由单调性必先经过 $(2,2)$ 再经过 $(3,3)$，路径数为
\[
\binom42\binom21\binom42=6\cdot2\cdot6=72.
\]
故
\[
|A\cup B|=120+120-72=168,\qquad 252-168=84.
\]
答案为 $84$。

\begin{pitfall}[单调性给出禁入点的天然顺序]
在单调格路上，若 $(a,b)$ 与 $(c,d)$ 满足 $a\le c$、$b\le d$，则任何经过这两点的路径必先经过 $(a,b)$。所以容斥里的交集不需要再分顺序——这是格路容斥比一般容斥简单的地方。但\textbf{若两点不可比较}（例如 $(2,3)$ 与 $(3,2)$），一条路径至多经过其中一个，交集为空。
\end{pitfall}

\chapter{组合构造与综合例题}
\begin{chapterintro}[章节导读]
计数回答“有多少”，构造回答“能不能造一个”。竞赛里的构造题往往不需要精巧的想法，只需要把条件逐条满足；真正难的是“证明你造的确实合法”。本章给出构造题的通用流程，并用不变量、图论计数与综合例题收束第三部分。
\end{chapterintro}

\section{构造题的标准流程}
\begin{enumerate}
\item \textbf{写下全部条件}：把题面转写成带量词的命题，明确“任意”与“存在”。
\item \textbf{找必要条件}：先推出一些“必须成立”的简单条件（奇偶性、计数和、度数），它们既是构造的指引，也是“为什么某些情形无解”的答案。
\item \textbf{按必要条件构造}：多数题在满足必要条件后就能构造；常用套路有：贪心、归纳（从 $n-1$ 到 $n$）、对称构造、分解成独立小块。
\item \textbf{验证合法性}：逐条核对条件，并检查边界（$n=1$、空集、相等的情形）。
\end{enumerate}

\begin{pitfall}[“看起来合法”不算证明]
构造题最常见的失分是“给了构造但没验证”。评分时要求：每一组输出都要能说明它满足每一条约束。第四部分博弈论第 2 章“下模仿棋”是一个范例——构造只是“镜像”，但必须论证“镜像永远合法”。
\end{pitfall}

\section{不变量}
\begin{definition}[不变量]
在操作过程中保持不变的量（或保持某种单调性的量）叫\textbf{不变量}。若能找到一个不变量在初始状态与目标状态取值不同，则目标不可达。
\end{definition}
常见不变量：奇偶性、模 $k$ 余数、总和、异或和、图的度数奇偶性。第二部分第 4 章“环上的固定步长”用的就是“编号模 $d$ 的余数不变”。

\begin{problem}{例：倒水问题}
有容量 $3$ 升与 $5$ 升的两个空壶，能量出 $4$ 升水吗？
\end{problem}
\hint 每次操作使某壶装满、倒空或倒入另一壶至满；不变量是“总水量”与“各壶水量都是整数”。

\sol 可以：$(0,0)\to(0,5)\to(3,2)\to(0,2)\to(2,0)\to(2,5)\to(3,4)$，此时 $5$ 升壶中剩 $4$ 升。一般结论：容量 $a,b$（互质）能量出 $1$ 到 $\max(a,b)$ 之间任意整数升——这正是裴蜀定理的构造性版本，第二部分第 1 章会给严格证明。

\section{图上的计数}
\begin{theorem}[握手引理]
任意无向图中，所有顶点的度数之和等于边数的两倍：
\[
\sum_v\deg(v)=2|E|.
\]
因此度数为奇数的顶点个数必为偶数。
\end{theorem}
\begin{proof}
每条边恰好为两个端点的度数各贡献 $1$，求和时每条边被数两次。
\end{proof}

\begin{theorem}[树的边数]
$n$ 个顶点的树恰有 $n-1$ 条边。
\end{theorem}
\begin{proof}[归纳]
$n=1$ 时 $0$ 条边。$n\ge2$ 时树必有叶子（度数为 $1$ 的顶点，否则沿边走会形成环或无限路径，与有限且无环矛盾）。删去一片叶子及其边，剩下一棵 $n-1$ 个顶点的树，由归纳假设有 $n-2$ 条边，加回删除的边得 $n-1$ 条。
\end{proof}

\begin{tip}[欧拉公式]
连通平面图满足 $V-E+F=2$（$F$ 含外部面）。它是平面图计数的基本工具，也能推出“平面简单图边数不超过 $3V-6$”。这部分内容在算法竞赛中出现频率不高，列在此处作为图论计数的入口。
\end{tip}

\section{综合例题}
\begin{problem}{综合 1：二项式系数求和}
计算 $\sum_{k=0}^{n}k\binom nk$ 与 $\sum_{k=0}^{n}\dfrac1{k+1}\binom nk$。
\end{problem}
\hint 第一问用求导；第二问积分，或注意到 $\dfrac1{k+1}\binom nk=\dfrac1{n+1}\dbinom{n+1}{k+1}$。

\sol 第一问 $n2^{n-1}$。第二问：
\[
\sum_{k=0}^{n}\frac1{k+1}\binom nk
=\frac1{n+1}\sum_{k=0}^{n}\binom{n+1}{k+1}
=\frac1{n+1}\left(2^{n+1}-1\right).
\]
最后一步用了 $\sum_{j}\dbinom{n+1}{j}=2^{n+1}$ 并去掉 $j=0$ 项。

\begin{problem}{综合 2：容斥 + 递推}
用 $1\times2$ 与 $2\times1$ 的骨牌完全覆盖 $2\times n$ 的棋盘，有多少种方法？
\end{problem}
\hint 按最左边一列怎么被覆盖分类。

\sol 设 $f_n$ 为覆盖 $2\times n$ 的方法数。最左一列要么被一个竖放的 $2\times1$ 覆盖（剩下 $2\times(n-1)$，$f_{n-1}$ 种），要么被两个横放的 $1\times2$ 覆盖（剩下 $2\times(n-2)$，$f_{n-2}$ 种）。故
\[
f_n=f_{n-1}+f_{n-2},\qquad f_0=1,\ f_1=1,
\]
即 $f_n=F_{n+1}$。

\begin{problem}{综合 3：抽屉原理与奇偶拆分}
证明：任意 $n+1$ 个不超过 $2n$ 的正整数中，必有两个数，其中一个整除另一个。
\end{problem}
\hint 把每个数写成“奇数部分 $\times$ $2$ 的幂”；$1$ 到 $2n$ 中只有 $n$ 个不同的奇数部分。

\sol 任取一个正整数 $a\le2n$，它可唯一写成
\[
a=2^{t}\cdot u,\qquad u\ \text{为奇数}.
\]
在 $1,\ldots,2n$ 中，奇数只有 $1,3,5,\ldots,2n-1$，共 $n$ 个，所以可能的 $u$ 只有 $n$ 种。现有 $n+1$ 个数，由抽屉原理，必有两个数 $a=2^{t_1}u$、$b=2^{t_2}u$ 具有相同的奇数部分 $u$。设 $t_1\le t_2$，则 $a\mid b$。得证。

这个证明是“\textbf{先找正确的分类，再套抽屉}”的典范：分类标准不是数的大小，而是\textbf{剥离 $2$ 的幂以后剩下的奇数核心}。

\begin{teach}[综合 3 讲评要点]
让学生先自己造抽屉。常见的错误尝试有：按“$\le n$ 与 $>n$”分两类（不奏效，因为两类内部没有整除关系）、按“是否偶数”分（同样不奏效）。等到他们发现这两条路都走不通，再给出奇数核心的分类。\textbf{抽屉造不出来，往往是因为分类标准选错了维度。}
\end{teach}

\begin{problem}{综合 4：计数与验证并重}
证明：任意 $n\ge2$ 个人中，若每两个人要么认识要么不认识，则“认识人数为奇数”的人数是偶数。
\end{problem}
\hint 握手引理。

\sol 由握手引理，$\sum_v\deg(v)=2|E|$ 是偶数。设度数为奇数的顶点集合为 $O$、度数为偶数的为 $E_v$，则
\[
\sum_{v\in O}\deg(v)=2|E|-\sum_{v\in E_v}\deg(v)
\]
是偶数减偶数，仍为偶数。而奇数个奇数相加为奇数、偶数个奇数相加为偶数，故 $|O|$ 必为偶数。

\begin{teach}[第三部分的结课检查]
让学生独立完成三道小题：
\begin{enumerate}
\item 证明 $\dbinom{2n}{n}$ 是偶数（$n\ge1$）——提示用 $\dbinom{2n}{n}=\dbinom{2n-1}{n-1}+\dbinom{2n-1}{n}$；
\item 求 $n$ 对括号的正确匹配数，并用反射原理从格路角度再证一次；
\item 判断命题“任意 $n+2$ 个不超过 $2n$ 的正整数中，必有两个数的差能被 $n$ 整除”的真假，并证明或给出反例。
\end{enumerate}
三题全对即可进入第四部分。第三题需要谨慎：先在 $n=2$、$n=3$ 上试探，确认命题真假再下笔。
\end{teach}
